\section{The calibrated velocity controller}
\label{sec:controller}

Given a velocity request $\bm v^*$, invert the model to get the effort that
\emph{should} produce it (feed-forward), then correct what the model got wrong
with feedback on the measured velocity $\bm v$:
\[
  \bm e = \underbrace{G^{-1}\left(\operatorname{clip}(\bm v^*) - \bm b\right)}_{\text{feed-forward}}
        + K_p\,\bm\epsilon + K_i\,\bm z,\qquad \bm\epsilon = \bm v^* - \bm v,
\]
with $\bm e$ finally clipped to $[-1,1]$. $\operatorname{clip}$ limits the
request to the measured velocity limits.

\paragraph{Bounded integral (anti-windup).} The integral state
\[
  z_i \leftarrow \operatorname{clip}\!\left(z_i + \epsilon_i\,\Delta t,\; -z_\text{max},\; z_\text{max}\right)
\]
removes steady errors such as a constant current the bias did not capture. The
clip matters: without it, a long period of saturated thrust (a request the boat
cannot meet) keeps accumulating error, and the boat overshoots badly once it
finally can respond. That is called \emph{windup}. The integral also resets when
commands stop.

\paragraph{Why feed-forward first.} With a good model the feed-forward term does
almost all of the work, and the feedback gains can stay small and stable. If
the model is poor, the controller-steps procedure (Section~\ref{sec:step}) shows
it as large steady-state error or overshoot.
\file{mhseals\_hardware/velocity\_controller.py}
